\documentclass[11pt]{article}

\usepackage[utf8]{inputenc}
\usepackage[T1]{fontenc}
\usepackage{lmodern}
\usepackage{amsmath}
\usepackage{amssymb}
\usepackage{amsthm}
\usepackage[hidelinks]{hyperref}
\usepackage{doi}

% Map fi/ffi/fl ligature glyphs to Unicode so PDF full-text extraction (Zenodo indexing) is correct.
\input{glyphtounicode}
\pdfgentounicode=1

\theoremstyle{definition}
\newtheorem{definition}{Definition}
\theoremstyle{remark}
\newtheorem{remark}{Remark}

\newcommand{\CV}{\operatorname{CV}}
\newcommand{\rl}{r_{\ell}}
\newcommand{\SPi}{S_{\Pi}}

\title{Ternary Phase Carrier: Colour Confinement, Generation Splitting,\\ and Koide Dispersion}

\author{J. Beau\\
Independent Researcher, France\\
\texttt{jerome.beau@cosmochrony.org}}

\date{2026}

\begin{document}

  \maketitle

  \begin{abstract}
    We record a single carrier hypothesis that would relate three otherwise separate facts of the fermionic
    sector: the exact non-observability of individual colours in the confined regime, the appearance of three
    massive generations, and the unit-dispersion (Koide) shape of the charged-lepton square-root masses.
    The proposal is one complex three-phase carrier $z_g = S + A\,e^{i(\delta + 2\pi g/3)}$, $g = 0,1,2$, read in
    two ways: its pure-phase part sums to zero over the three cube-root phases ($1 + \omega + \omega^2 = 0$), a
    collective and exact cancellation identified with a colour-neutral class; its real part $S + A\cos(\delta +
    2\pi g/3)$ is the square-root mass profile of the companion note~\cite{Beau2026kud}.
    In this reading the singlet amplitude $S$ is the master dial: $S = 0$ is the pure-phase, sum-zero regime
    (colour, confined); $S \neq 0$ opens the real-offset regime (positive masses); the balanced point $S =
    \lVert\text{doublet}\rVert$ is singlet--doublet equipartition (Koide, $\CV = 1$); and the phase $\delta$ places
    one generation near a node of the profile (the hierarchy).
    This is a \emph{speculative bridge}, not a derivation.
    It does not derive the Standard-Model charge assignments, anomaly cancellation, or the charged-lepton masses.
    It identifies a candidate common origin for the colour triplet and the generation triplet, and states the
    explicit construction charge sheet that any realisation would have to discharge.
  \end{abstract}

  \noindent\textbf{Keywords:}
  ternary phase carrier; colour confinement; generation splitting; Koide relation; unit dispersion; sum-zero
  cancellation; projective regime; singlet--doublet equipartition; speculative bridge.

  \newpage

  \tableofcontents

  \newpage

  \section{Position and status}
    \label{sec:position}

    The Cosmochrony fermionic sub-programme carries three facts that are, in the Standard Model, logically
    independent.
    First, colour is an exact, unbroken structure: the three colours are degenerate and no single colour is
    extractable as an observable, a non-observability the framework already reads as the non-injective colour
    projection of ENI~\cite{Beau2026l} and tests as an exactly degenerate colour triplet on
    $\mathrm{Heis}_3(\mathbb{Z}/q\mathbb{Z})$~\cite{Beau2026a35, Beau2026a36}.
    Second, there are three massive generations, hierarchically split.
    Third, the charged-lepton square-root masses have unit relative dispersion, the scale-free content of the
    Koide relation~\cite{Beau2026kud}.
    The number three recurs in all three, and their coincidence is not explained by treating colour and generation
    as independent labels.

    This note records the sharpest form of a hypothesis that would tie the three together with a single object,
    without multiplying mechanisms.
    Its status is deliberately bounded, and stated here once so that nothing downstream is mistaken for a result.

    \begin{quote}
      This note proposes a unifying carrier hypothesis.
      It does not derive the Standard-Model charge assignments, anomaly cancellation, or the charged-lepton
      masses.
      It identifies a candidate bridge between the colour triplet and the generation triplet.
    \end{quote}

    The reason for a separate note is exactly this bounded status.
    The companion unit-dispersion note~\cite{Beau2026kud} is a reliable target-invariant statement about the
    charged leptons; the present hypothesis is broader~---~it relates colour, generation, and mass~---~and
    speculative, so it is kept apart in order not to contaminate the clean results.

  \section{The ternary phase carrier}
    \label{sec:carrier}

    \begin{definition}[Ternary phase carrier]
      \label{def:carrier}
      Let $S \ge 0$ (the singlet, or democratic, amplitude), $A \ge 0$ (the doublet, or oscillating, amplitude),
      and $\delta$ a phase.
      The ternary phase carrier is the triple of complex numbers
      \begin{equation}
        \label{eq:carrier}
        z_g = S + A\,e^{i\left(\delta + \frac{2\pi g}{3}\right)}, \qquad g = 0, 1, 2 .
      \end{equation}
    \end{definition}

    The carrier admits two readings, which are the two halves of one object rather than two hypotheses.

    \paragraph{Pure-phase reading (sum-zero).}
      The oscillating part alone, summed over the three equally spaced phases, cancels exactly:
      \begin{equation}
        \label{eq:sumzero}
        \sum_{g=0}^{2} e^{i\left(\delta + \frac{2\pi g}{3}\right)}
        = e^{i\delta}\,(1 + \omega + \omega^2) = 0, \qquad \omega = e^{2\pi i/3} .
      \end{equation}
      This is not the absence of the three components but their collective annulment: the three phases are present
      at the level of the internal structure, yet no single phase is extractable, because the observable is their
      vanishing sum.
      This is precisely the language of a colour-neutral class~---~the colours are internally present but
      projectively invisible, the sum falling into the neutral sector~\cite{Beau2026l, Beau2026a35, Beau2026a36}.

    \paragraph{Real reading (offset profile).}
      The real part of~\eqref{eq:carrier} is
      \begin{equation}
        \label{eq:realprofile}
        \Re z_g = S + A\cos\!\left(\delta + \frac{2\pi g}{3}\right),
      \end{equation}
      which is exactly the cyclic square-root mass profile of the companion note~\cite{Beau2026kud} (there written
      with $S = \bar r$ and $A = \sqrt2\,\bar r$).
      The democratic offset $S$ makes the three values real and positive; the oscillating part supplies their
      spread; the phase $\delta$ fixes which generation is light.

    The complex sum-zero of the first reading and the real near-node of the second are not the same operation~---~one
    is an exact, collective, complex cancellation, the other an approximate, individual, real amplitude minimum.
    What relates them is precisely the offset $S$: suppressing $S$ leaves the pure-phase regime whose sum
    vanishes; restoring $S$ opens the real-positive regime in which masses appear.
    A single carrier, two lectures, with $S$ the switch between them.

  \section{The singlet amplitude as the master dial}
    \label{sec:dial}

    Reading~\eqref{eq:carrier} through the amplitude $S$ organises the three facts as three values of one dial
    together with one phase.

    \begin{center}
      \small
      \begin{tabular}{@{}l l p{0.34\textwidth}@{}}
        \hline
        regime & carrier reading & phenomenon \\
        \hline
        $S = 0$ & pure phase, $\sum_g e^{i\phi_g} = 0$ & colour-neutral (confined) \\
        $S \neq 0$ & real offset, positive values & massive sector opens \\
        $S = \lVert\text{doublet}\rVert$ & singlet--doublet equipartition & Koide dispersion ($\CV = 1$) \\
        $\delta$ near a node & one real amplitude annulled & hierarchy (light electron) \\
        \hline
      \end{tabular}
    \end{center}

    The equipartition row is the point of contact with the companion note.
    For three cosines $120^\circ$ apart, $\langle \cos^2 \rangle = 1/2$, so the oscillating component has effective
    norm $A/\sqrt2$; equality of the differentiating and democratic norms, $A/\sqrt2 = S$, is the single condition
    \begin{equation}
      \label{eq:equipartition}
      \frac{A}{\sqrt2} = S
      \quad\Longleftrightarrow\quad
      \CV(\rl) = 1
      \quad\Longleftrightarrow\quad
      \theta(\rl, \mathbf 1) = 45^\circ,
    \end{equation}
    the unit-dispersion form of Koide~\cite{Beau2026kud}, with $\theta$ the angle to the democratic axis
    $\mathbf 1 = (1,1,1)$.
    The amplitude $\sqrt2$, the Koide angle $45^\circ$, and $\CV = 1$ are one content~---~singlet--doublet
    equipartition of the carrier~---~and in the present language this is the balanced setting of the dial $S$.

    The whole hypothesis then compresses to one sentence.

    \begin{quote}
      Colour is the pure-phase, sum-zero regime of the ternary carrier; mass is the real-offset, singlet-balanced
      regime of the same carrier; hierarchy is the residual phase placement relative to a node.
    \end{quote}

    \begin{remark}[What the compression does and does not buy]
      \label{rem:compress}
      The gain is economy: one carrier and one dial replace three separate postulates (a colour structure, a
      generation replication, a dispersion coincidence).
      The cost is unchanged.
      The dial reading fixes neither the balanced value $A = \sqrt2\,S$ (the form lock) nor the node phase $\delta$
      (the placement lock); both remain open exactly as in the companion note~\cite{Beau2026kud}, and the placement
      lock stays the harder, exponentially sensitive residue.
      The carrier reorganises the open questions; it does not close them.
      In particular the dial is a \emph{reparametrisation} of the carrier's regimes, not a derived dynamics: the
      $S = 0 \to S \neq 0$ passage that ``opens the massive sector'' is constrained, and its most natural
      realisation is already closed (Section~\ref{sec:chargesheet}, item~5).
    \end{remark}

    \begin{remark}[Equipartition as a marginal-stability edge, not a tuned value]
      \label{rem:edge}
      The balanced setting $A = \sqrt2\,S$ need not be read as a value the projection is tuned to, but as the one
      configuration that survives a viability filter.
      The two flanking regimes are sterile: as $\CV \to 0$ the three components coincide and no generation is
      discernible (uniformity), while as $\CV \to \sqrt2$ one component dominates and the other two are annulled (no
      plurality); only the midpoint $\CV = 1$ carries at once a common structure and a genuine multiplicity, the sole
      regime in which a three-generation spectrum can both exist and propagate.
      Equipartition would then be what remains once the uniform and the dominated regimes have failed to be
      prolific~---~an edge of viability rather than a coincidence.
      Two things make this more than an image.
      First, a stability boundary is an \emph{equality} (marginal means exactly on the edge), not an interval, so it
      can fix a sharp value; this is the mode in which the programme's admissibility and saturation conditions
      already operate~\cite{Beau2026bim, Beau2026l}.
      Second, and conversely, only a boundary can be fundamental at all: an interior point of a viable band would be
      contingent (why this value and not a neighbour?), whereas an edge is forced, so being an edge is not incidental
      to the hypothesis but is what would qualify $\CV = 1$ as fundamental.
      In this reading the two locks of Remark~\ref{rem:compress} unify: equipartition is the viability
      \emph{constraint} (stay on the edge), and the node placement of the phase $\delta$ is the \emph{objective}
      (differentiate as far as possible), the observed spectrum pushing one amplitude near annulment while remaining
      exactly on the edge $\CV = 1$.
      In canonical form: the balanced value $\CV = 1$ should be read, if fundamental at all, not as a tuned interior
      point of a viable interval but as a marginal-stability edge~---~the exact boundary at which generation
      differentiation becomes maximal without destroying the common carrier.
      This is speculative.
      A derivation would require exhibiting the actual filter~---~a relaxation or an admissibility saturation~---~whose
      unique marginal survivor is $\CV = 1$, and showing that the two flanks are genuinely inadmissible in the
      corpus's own terms, not merely intuitively sterile.
      Crucially, the edge reading is \emph{not} an independent argument for $\CV = 1$.
      The candidate filter of Section~\ref{sec:filter} \emph{defines} $\CV \le 1$ as the admissible band, so that
      $\CV = 1$ is its boundary holds by construction; the marginal-stability language is a reformulation of the
      filter, not a second and independent piece of evidence.
      All the content lies in justifying the filter itself: absent that, this remark restates the hypothesis and
      must not be mistaken for a physical derivation of it.
    \end{remark}

  \section{A candidate admissibility filter}
    \label{sec:filter}

    The carrier reading suggests a natural but still hypothetical admissibility filter on the generation sector.
    Decompose the real square-root profile into its democratic and differentiating parts,
    \begin{equation}
      \label{eq:decomp}
      r = r_{\mathrm{sing}} + r_{\mathrm{dbl}},
      \qquad r_{\mathrm{sing}} = \bar r\,\mathbf 1,
      \qquad \langle r_{\mathrm{dbl}}, \mathbf 1 \rangle = 0 .
    \end{equation}
    Then $\lVert r_{\mathrm{sing}} \rVert = \sqrt3\,\bar r$ and $\lVert r_{\mathrm{dbl}} \rVert = \sqrt3\,\sigma_r$, so
    \begin{equation}
      \label{eq:ratio}
      \frac{\lVert r_{\mathrm{dbl}} \rVert}{\lVert r_{\mathrm{sing}} \rVert} = \frac{\sigma_r}{\bar r} = \CV(r),
    \end{equation}
    and the Koide value is exactly the balanced condition
    \begin{equation}
      \label{eq:balanced}
      \CV(r) = 1 \quad\Longleftrightarrow\quad \lVert r_{\mathrm{dbl}} \rVert = \lVert r_{\mathrm{sing}} \rVert .
    \end{equation}

    This equality should \emph{not} be read as a positivity edge of the real profile, because positivity gives the
    wrong boundary.
    For the continuous envelope $S + A\cos\phi$, positivity requires $A \le S$, hence
    \begin{equation}
      \label{eq:envelope}
      \CV = \frac{A}{\sqrt2\,S} \le \frac{1}{\sqrt2} .
    \end{equation}
    This is precisely the obstruction met by the cascade ladder in the companion unit-dispersion
    note~\cite{Beau2026kud}: the continuous envelope cannot reach the Koide value.
    At the opposite end, the full positive cone of three discrete components has boundary $\CV = \sqrt2$, where one
    component vanishes.
    The Koide value $\CV = 1$ lies strictly between these two positivity boundaries.\footnote{%
      Checked numerically: strictly positive triples with $1 < \CV < \sqrt2$ are abundant (e.g.\ $(0.11, 34.30,
      1.34)$ gives $\CV = 1.33$ with all three components positive), so $\CV > 1$ is not excluded by positivity.
      The continuous-envelope bound $\CV \le 1/\sqrt2$ and the positive-cone boundary $\CV = \sqrt2$ are the only
      positivity edges, and $\CV = 1$ lies strictly between them.}

    The candidate filter is therefore not positivity, nor a naive Born--Infeld saturation of the real envelope.
    It is a discrete capacity balance,
    \begin{equation}
      \label{eq:filter}
      \boxed{\ \lVert r_{\mathrm{dbl}} \rVert \le \lVert r_{\mathrm{sing}} \rVert\ } .
    \end{equation}
    The differentiating generation doublet may spend at most the capacity carried by the democratic singlet.
    Projection locking then selects the saturated admissible edge,
    \begin{equation}
      \label{eq:filter-sat}
      \boxed{\ \lVert r_{\mathrm{dbl}} \rVert = \lVert r_{\mathrm{sing}} \rVert\ }
      \quad\Longleftrightarrow\quad \CV(r) = 1 .
    \end{equation}

    In this corrected reading the two flanks have different statuses.
    For $\CV < 1$, the generation doublet has unused capacity: the sector is under-resolved and not maximally
    differentiating.
    For $\CV > 1$, the doublet exceeds the singlet capacity: the differentiating sector would carry more projective
    weight than the common carrier.
    The open question is not whether such triples can be positive~---~they can~---~but which projective capacity or
    entropy principle forbids this excess.

    Two distinct notions of discreteness must be separated here, and separating them reconciles this filter with the
    companion note.
    The discrete-to-continuous bridge that the companion note~\cite{Beau2026kud} calls its deepest lock is specific
    to its \emph{max-entropy} route, where $\CV = 1$ is read off a continuous probability density and does not
    transfer to three discrete masses.
    In the present carrier framing that transfer does not arise: three phases exactly $120^\circ$ apart inherit the
    amplitude-fixed dispersion \emph{exactly}, $\CV = A/(\sqrt2\,S)$ independently of the placement $\delta$, so
    $\CV = 1 \iff A = \sqrt2\,S$ holds directly on the discrete triple.
    What is genuinely discrete here is instead a positivity fact about the carrier envelope: the continuous profile
    $S + A\cos\phi$ is nonnegative only up to $\CV = 1/\sqrt2$, whereas $\CV = 1$ (with $A = \sqrt2\,S$) has the
    envelope negative between the sampled points, so the Koide value is admissible only because the carrier is read
    at three points and not as a whole envelope.
    The remaining question is therefore not a density-to-discrete transfer but the single amplitude condition
    $A = \sqrt2\,S$~---~the saturated capacity balance~---~meaningful only at the generation-discrete level.

    \begin{remark}[Status of the filter]
      \label{rem:filter-status}
      The filter above is a hypothesis, not a theorem.
      Its value is to locate the remaining mechanism sharply: the positivity and continuous-envelope routes give the
      wrong boundary and are therefore excluded as sources of $\CV = 1$.
      A derivation would have to produce the capacity inequality
      $\lVert r_{\mathrm{dbl}} \rVert \le \lVert r_{\mathrm{sing}} \rVert$ from an existing projective
      quantity~---~the projection entropy $\SPi$~\cite{Beau2026l}, a generation-level capacity functional, or an
      admissibility principle on the ternary carrier.
      Until such an object is exhibited, the filter remains a structured open test.
    \end{remark}

  \section{The construction charge sheet}
    \label{sec:chargesheet}

    A carrier hypothesis of this kind is a front of construction, not a veto and not a result.
    Its cautious form must be maintained: the leptons are \emph{not} to be read as hidden coloured composites,
    which would be experimentally dangerous, since leptons are observed as exact colour singlets and as pointlike.
    The safe statement is that the lepton sector retains a projective trace of the ternary carrier while its colour
    component is exactly neutralised in the observable~---~the erased content is unobservable but may leave a
    residual invariant, which is a distinction the projective framework is built to express~\cite{Beau2026l}.
    Any realisation of the hypothesis must then discharge, at least, the following.
    \begin{enumerate}
      \item Produce an \emph{exact} colour singlet (the sum-zero regime, not an approximate one).
      \item Keep three \emph{discernible} generations after neutralisation, so that colour neutrality does not
        collapse the generation count (the colourless leptons must still carry three generations).
      \item Reproduce the leptons' weak isospin and hypercharge assignments, not merely their colourlessness.
      \item Preserve gauge-anomaly cancellation generation by generation, as in the Standard Model.
      \item Explain the selection rule: why the quark regime keeps an observable colour residue while the lepton
        regime cancels it exactly.
    \end{enumerate}
    Items~1--4 are concrete demands rather than obstructions of principle.
    Item~5, however, is sharper than concrete work still to do, because its most natural realisation is already
    closed by the corpus.
    The obvious reading of the $S = 0 \to S \neq 0$ passage would be that neutralising the colour triplet is itself
    what produces the differentiated generation spectrum~---~turning the dial $S$ up along the colour route.
    That route is excluded.
    The colour triplet is tested \emph{exactly degenerate} on $\mathrm{Heis}_3(\mathbb{Z}/q\mathbb{Z})$~---~rank one,
    equal singular values, and vanishing sum (a pure doublet, no singlet component)~\cite{Beau2026a35, Beau2026a36};
    and the projection is forgetful, with strictly positive projection entropy $\SPi > 0$~\cite{Beau2026l}, so it
    cannot manufacture a spread from a rank-one, structureless input.
    A lossy map does not create the differentiation it is not given: the generation splitting therefore cannot be
    read off the colour neutralisation itself, and the breaking that carries $S \neq 0$ must be seeded from outside
    the colour sector.
    The dial language must not obscure this~---~the amplitude $S$ is a reparametrisation of the carrier's two
    regimes, not a derived dynamics interpolating between them, and its most natural dynamical realisation, colour
    turning into flavour under the projection, is precisely the one the corpus has closed.
    The quark--lepton relation is not exotic in isolation~---~Pati--Salam unification already treats lepton number
    as a fourth colour, and is bounded (by proton-decay and rare-decay limits) rather than excluded; but the present
    hypothesis is a further and distinct step, relating colour to the \emph{generation} triad, and for that step the
    colour-neutralisation route to the splitting is not merely unbuilt but closed.
    Items~4 and~5 are where the real work sits, and item~5 must look beyond colour neutralisation for the source of
    the breaking.

  \section{Status}
    \label{sec:status}

    The content of this note is a hypothesis with a precise shape and a bounded claim.
    The carrier $z_g = S + A\,e^{i(\delta + 2\pi g/3)}$ has an exact pure-phase sum-zero~\eqref{eq:sumzero} and a
    real part~\eqref{eq:realprofile} that coincides with the square-root mass profile of the companion
    note~\cite{Beau2026kud}; these two statements are exact algebra.
    The organisation of colour, mass, and Koide dispersion as three settings of the singlet dial $S$ plus one
    phase $\delta$ is a \emph{reading}, internally consistent and economical, but not derived: the corpus supplies
    the colour sum-zero~\cite{Beau2026l, Beau2026a35, Beau2026a36} and the real square-root profile~\cite{Beau2026kud}
    as separate facts, and their unification as the imaginary-phase and real-offset faces of one carrier is the
    hypothesis, not a theorem.
    The charge sheet of Section~\ref{sec:chargesheet} is the price of turning the reading into a mechanism, and it
    is unpaid.
    In particular this note does not fix the equipartition value $A = \sqrt2\,S$, does not fix the node phase
    $\delta$, and does not derive the Standard-Model quantum numbers or anomaly cancellation.
    Whether the projection entropy $\SPi$ of the non-injective map~\cite{Beau2026l}, or a related variational
    principle, forces the balanced setting of the dial is the same open test recorded for the unit-dispersion
    target~\cite{Beau2026kud}, now phrased on the carrier.
    The working conclusion is only that the three facts \emph{can} be written as one carrier read three ways~---~a
    candidate bridge, held at the level of a speculative outlook.

  \newpage

  \bibliographystyle{plain}
  \bibliography{cosmochrony-bibliography}

\end{document}
